\section{Support Vector Machines. Метод на опорните вектори.}	
\includegraphics[scale=0.5]{svn.png} \\
Разделяща хиперравнина. \\
Recall: Хиперравнина се задава с $\beta_0 + \beta_1 X_1 + \dots + \beta_p X_p = 0$. Нормален вектор на тази хиперравнина
е $(\beta_1, \dots \beta_p)$ и той задава положителната посока. Уравнението може да се запише и като $\beta_0 + (\boldsymbol \beta, \mathbf X) = 0 $, където 
$\boldsymbol \beta = (\beta_1, \dots \beta_p) $ и $\mathbf X = (X_1, \dots, X_p)$. Една хиперравнина разделя пространството на две части, една с $\beta_0 + (\boldsymbol \beta, \mathbf X) > 0 $, и една с $\beta_0 + (\boldsymbol \beta, \mathbf X) < 0 $. Ако имаме точки в пространството, които можем да разделим с хиперравнина, и дадем labels на тези точки $y_i = 1$ и $y_i=-1$, то 
$$y_i(\beta_0 + \beta_1 x_{i1} + \dots \beta_p x_{ip}) > 0 .$$
Оптимизационна задача:
\begin{align}
%\begin{equation}
&\text{да се максимизира } M \\ 
%\end{equation}
%\begin{equation}
&\text{така че } \sum_{j=1}^{p} \beta_j^2 = 1 \\
%\end{equation}
%\begin{equation}
&	y_i(\beta_0 + \beta_1 x_{i1} + \dots \beta_p x_{ip}) \geq M \text{ за } i=1,\dots,n
%\end{equation}
\end{align}
Втора оптимизационна задача(Support Vector Classifier):
\begin{align}
&\text{да се максимизира } M \\ 
&\text{така че } \sum_{j=1}^{p} \beta_j^2 = 1 \\
&	y_i(\beta_0 + \beta_1 x_{i1} + \dots \beta_p x_{ip}) \geq M(1-\epsilon_i) \text{ за } i=1,\dots,n \\
& \text{ и } \epsilon_i \geq 0, \sum_i e_i = C 
\end{align}

Support Vector Machines:\\
Решението на втора оптимизационна задача може да се запише във вида
 $$ f(x) = \beta_0 + \sum_{i=1}^n \alpha_i \langle x,x_i \rangle , $$
 и тогава на нас са ни достатъчни скаларните произведения $ \langle x_i,x'_j \rangle $ за всички $ \binom{n}{2} $ двойки $x_i,x'_j $, $i,j = 1, \dots,n $, за да пресметнем модела.	
Вместо 	$ \langle x_i,x'_j \rangle $, може да вземем произволна kernel фунцкия $ K( x_i,x'_j ) $. Примерът със скаларното произведение 
е с линейна фунцкия, но може и нелинейна. \\

Упражнения: Conceptual 1,3 


Solution
Main	  advantages of	  SVMs:
\begin{itemize}
\item Has rigorous theoretical foundation
\item Performs classification more accurately than most other methods in applications, especially for high dimensional data
\item They can applied also to non linear classification problems by using kernel functions. The $"$kernel trick$"$ allows that these problems computationally tractable
\item  Different kernels can be plugged into the same learning machinery and studied independently of it.
\end{itemize}

Math limitations of the SVM:
\begin{itemize}
	\item Works only in a real-valued space.
	\item For a categorical attribute, we need to convert its categorical values to numeric values
	\item Does only two-class classification
	\item For multi-class problems, some strategies can be applied e.g. one-against-rest
	\item The hpyerpalne produced by SVM is hard to understand
\end{itemize}
\newpage
